% Reconstructed from the complete audited author transcriptions after CP09.
% Root, B and D passed the full mathematical transcription and the stated
% primary-source checks. This is a reconstruction, not recovered old bytes.
\begin{theorem}[Proved: A local rank-linear bound in characteristic $p$]
\label{D:local-p-linear-H1}
Let $(R,\mathfrak m,k)$ be a complete excellent normal local ring of
dimension two, of characteristic $p>0$, with algebraically closed residue
field. Put $W=\operatorname{Spec}R\setminus\{\mathfrak m\}$. There is a
constant $c_R$ such that every locally constant sheaf $M$ of rank $r$
over any finite field of characteristic $p$ satisfies
\[
 \dim H^1(W,M)\le c_Rr.
\]
The same conclusion holds for the puncture of an excellent strictly
henselian normal local surface ring with algebraically closed residue
field, with the constant supplied by its completion.
\end{theorem}
\begin{proof}
We first use $\mathbf F_p$ coefficients. The bundle
$\mathcal O_W\otimes_{\mathbf F_p}M$ has a coherent reflexive extension
$E$ to $R$. It is maximal Cohen--Macaulay. Its Frobenius linearization
$F_R^*E\to E$ is an isomorphism on $W$, and its kernel and cokernel
are supported at $\mathfrak m$. Set $H=H^2_{\mathfrak m}(E)$. The top
Cech cokernel for two parameters commutes with every right-exact tensor
functor. Consequently the induced structural map
\[
 R\otimes_{F,R}H\longrightarrow H
\]
is an isomorphism. This assertion does not require flatness of Frobenius:
the finite-length defects of the original linearization have zero top
local cohomology.

Let $H_{\rm nil}$ be the submodule of elements killed by some iterate of
the semilinear Frobenius. The structural isomorphisms show that, for
$z\in H_{\rm nil}$, the element $1\otimes z$ vanishes in every
sufficiently high Frobenius tensor power. Thus $H_{\rm nil}\subset
0_H^*$, with test multiplier $1$. Write $I$ for the injective hull of
$k$. By local duality $H$ embeds in a finite direct sum of copies of
$I$. Persistence of tight closure carries $H_{\rm nil}$ into the
tight closure of zero in this direct sum. The ideal
\[
 J=\operatorname{Ann}_R(0_I^*)
\]
therefore annihilates $H_{\rm nil}$. It is $\mathfrak m$-primary or is
$R$: the non-strongly-$F$-regular locus is contained in the closed point,
since every other local ring is regular. We use precisely the
injective-hull test ideal and its localization/support assertions in
\citet[Theorem 7.1(1), (2), (5)]{C_LyubeznikSmith2001}.

Let $D(-)=\operatorname{Hom}_R(-,I)$ be Matlis duality and set
\[
 P=D(H)=\operatorname{Hom}_R(E,\omega_R),\qquad
 N=D(H/H_{\rm nil}).
\]
Then $P$ is maximal Cohen--Macaulay of rank $r$, and
$JP\subset N\subset P$. Choose once a parameter ideal
$Q\subset\mathfrak mJ$. Since $\mathfrak mJP\subset\mathfrak mN$,
\[
 \dim_k N/\mathfrak mN
 \le \operatorname{length}(P/\mathfrak mJP)
 \le \operatorname{length}(P/QP)
 =r\operatorname{length}(R/Q).
\]
The last equality is the multiplicity formula for a parameter ideal on
a maximal Cohen--Macaulay module over a domain.

Frobenius is injective on $H/H_{\rm nil}$. A Frobenius-fixed element
$z$ of this Artinian module lies in its socle: for $a\in\mathfrak m$,
we have $a^{p^e}z=0$ for large $e$, so injectivity applied to
$F^e(az)$ gives $az=0$. Moreover $\mathbf F_p$-independent fixed vectors
are $k$-independent. Indeed, in a minimal $k$-linear relation we normalize
one coefficient to $1$, apply Frobenius, and subtract; minimality forces
all coefficients to lie in $\mathbf F_p$. It follows that
\[
 \dim_{\mathbf F_p}(H/H_{\rm nil})^{F=1}
 \le \dim_k\operatorname{Soc}(H/H_{\rm nil})
 =\dim_k N/\mathfrak mN\le c_Rr.
\]
A fixed vector in $H_{\rm nil}$ is zero, so the same bound holds for
$H^{F=1}$.

We spell out the Artin--Schreier comparison for a possibly nonfree
coherent module. A semilinear endomorphism of a finite module can be
lifted, after a finite free presentation, to a matrix of additive
polynomials $Cz^{[p]}-z$. Its differential is $-I$. Over an algebraically
closed residue field this additive polynomial map is surjective: its
image is an open subgroup of a connected vector group. Hensel lifting
then solves the equations over a strict local ring. Thus $F-1$ is
etale-locally surjective on $E$, and it is surjective on $E(R)$ here.
On $W$ its kernel is $M$. Since $H^1(W,\widetilde E)=H$, the resulting
long exact sequence embeds $H^1(W,M)$ into $H^{F=1}$. This proves the
assertion over $\mathbf F_p$. Restriction of scalars from
$\mathbf F_{p^f}$ multiplies both the rank and the cohomology dimension
by $f$, which cancels.

Finally let $R_0$ be the indicated strictly henselian ring. A connected
finite etale cover of its puncture normalizes to a finite normal domain
$B$ over $R_0$. Henselianity makes $B$ local. Excellence makes its
completion normal, and that completion is a local domain. Hence the
cover remains connected on the completed puncture. We obtain a
surjection of fundamental groups from the completed puncture to the
original one, and inflation gives the required injection on $H^1$ with
every finite local system. This proves the final assertion.
\end{proof}

\begin{theorem}[Proved: Uniform coherent cohomology for unit reflexive sheaves]
\label{D:unit-reflexive-linear-cohomology}
Let $S$ be a normal proper surface over an algebraically closed field of
characteristic $p>0$. There is a constant $c_S$ such that every reflexive
sheaf $E$ of rank $r$ equipped with an isomorphism
$F_S^{[*]}E\simeq E$ satisfies
\[
 h^0(S,E)\le r,\qquad h^1(S,E)\le c_Sr,\qquad
 h^2(S,E)\le c_Sr.
\]
Here $F^{[*]}$ denotes the reflexive hull of Frobenius pullback. If $E$
is obtained from a finite $\mathbf F_p$ local system on the regular
locus, then its fixed-section sheaf $L$ satisfies
$\dim_{\mathbf F_p}H^1(S,L)\le c_Sr$.
\end{theorem}
\begin{proof}
Frobenius scaling of the numerical first Chern class and second Chern
character gives $c_1(E)=0$ numerically and $\operatorname{ch}_2(E)=0$
\citep[Theorem 4.5(2)]{C_Langer2025}. The rank-linear correction in
Riemann--Roch for proper normal surfaces therefore gives
\[
 |\chi(S,E)-r\chi(S,\mathcal O_S)|\le c_0r
\]
\citep[Theorem 0.1]{C_Langer2026}. Properness, rather than projectivity
of $S$, is the hypothesis used here.

We fix a projective resolution $\pi:Y\to S$ with $Y$ smooth and
projective, and put $F=\pi^{[*]}E$. There is a fixed coherent ideal on
$S$, supported at its singular points, which annihilates the stable
endomorphism groups of all local maximal Cohen--Macaulay modules.
Equivalently, multiplication by each of its elements factors through a
finite free module. We use
\citet[Theorem 1.1 and Remark 2.5]{C_Kimura2025}, or the precise
factorization in \citet[Lemma 1.8]{C_Langer2026}.
If $z\operatorname{id}_E$ factors through a free module, Frobenius
pullback shows that $z^p$ annihilates the cokernel of
$F_S^*E\to F_S^{[*]}E$: take the reflexive hull of the factorization,
and use that the intervening free module already equals its hull.
After principalizing the fixed ideal of these $p$th powers, we obtain
a fixed effective divisor $A$ on $Y$ and injections
\[
 F(-A)\longrightarrow F_Y^*F\longrightarrow F.
\]
No constant in this sandwich depends on $E$ or $r$.

Fix an ample divisor $H$ and a smooth very ample curve $C\subset Y$
not contained in any of the fixed determinant-zero divisors. The maps
remain generically injective, hence injective as bundle maps, on $C$.
We do not require $C$ to be disjoint from $A$. Pulling back a maximal
slope subbundle by Frobenius gives, both on $Y$ and on $C$,
\[
 p\mu_{\max}(F)\le\mu_{\max}(F),\qquad
 p\mu_{\max}(F^\vee)\le\mu_{\max}(F^\vee)+A\cdot H,
\]
with $A\cdot C$ in the curve version. Thus the maximal slope of $F$
is nonpositive and that of $F^\vee$ is bounded by a fixed constant.
We have $h^0(F)\le r$. One direct verification takes the subsheaf
generated by all global sections: its determinant has nonnegative
degree, while the slope bound makes that degree nonpositive. Independent
generating sections therefore have determinant with no divisorial zero;
properness then makes every global section a constant linear combination
of at most $r$ of them.

Let $W=\pi^{[*]}\omega_S$, a fixed line bundle. Duality and restriction
give
\[
 H^2(S,E)^\vee=\operatorname{Hom}_S(E,\omega_S)
 \hookrightarrow H^0(Y,F^\vee\otimes W).
\]
Choose a fixed $m$ so that $F^\vee\otimes W(-mC)$ has negative maximal
slope. Successive restrictions to $C$ reduce the last $H^0$ to $m$
curve bundles whose maximal slopes are bounded by a fixed constant.
For a bundle $G$ on $C$ with maximal slope at most $a$, choose a fixed
effective divisor $D$ of degree greater than $a$; then
$H^0(C,G(-D))=0$ and $h^0(C,G)\le r\deg D$. We obtain
$h^2(S,E)\le c_1r$. Riemann--Roch now gives $h^1(S,E)\le c_Sr$.

The finite-free-presentation and Hensel argument in
Theorem~\ref{D:local-p-linear-H1} gives an exact etale sequence
\[
 0\longrightarrow L\longrightarrow E
 \xrightarrow{F-1}E\longrightarrow0.
\]
On the finite-dimensional space $H^0(S,E)$ the additive polynomial
$F-1$ is likewise surjective. Hence $H^1(S,L)$ injects into
$H^1(S,E)^{F=1}$. Independent fixed vectors over $\mathbf F_p$ are
independent over $k$, by the minimal-relation argument already used.
This proves the last assertion.
\end{proof}

\begin{theorem}[Proved: All-coefficient linear bounds on punctured proper surfaces]
\label{D:punctured-surface-all-characteristics-H1}
Let $S$ be a proper scheme of dimension at most two over an algebraically
closed field of characteristic $p>0$. Let $T$ be a finite closed subset
having codimension at least two in every irreducible component which it
meets, and let $U$ be a connected component of $S\setminus T$. There is
a constant $c_{S,T}$ such that, for every finite field $\Lambda$ and
every finite-dimensional continuous $\Lambda[\pi_1(U)]$-module $M$,
\[
 \dim_\Lambda H^1(\pi_1(U),M)\le c_{S,T}\dim_\Lambda M.
\]
For an open subgroup $H\subset\pi_1(U)$ the bound is
$c_{S,T}[\pi_1(U):H]\dim M$. In particular the maximal prosolvable
quotient of $\pi_1(U)$ is finitely generated. Full finite generation is
still subject to the nonabelian-crown condition of
Proposition~\ref{D:cohomology-crown-criterion}.
\end{theorem}
\begin{proof}
First suppose $S$ is a normal surface and use $\mathbf F_p$ coefficients.
For $j:U\hookrightarrow S$, the canonical reflexive extension of
$\mathcal O_U\otimes M$ is unit. Theorem~\ref{D:unit-reflexive-linear-cohomology}
bounds $H^1(S,j_*M)$ by a constant times the rank. The low-degree Leray
sequence has remaining term $H^0(S,R^1j_*M)$, supported at the finitely
many omitted points. Its stalks are bounded by
Theorem~\ref{D:local-p-linear-H1}; the completion comparison in that
theorem applies to each strict henselian stalk. This proves the
characteristic-$p$ bound, and restriction of scalars proves it over
every finite field of that characteristic. For coefficient characteristic
different from $p$ we apply
Proposition~\ref{D:surface-curvewise-conductor}, including its
rank-independent conductor and rank-linear Betti bounds.

For general $S$, reduction does not change its etale site, and the
normalization $\nu:X\to S_{\rm red}$ is finite. On each normalized
surface component we use the normal case. The components of dimension
at most one are proper and have finitely generated fundamental groups,
so their $H^1$ also has a fixed rank-linear bound. For the finite
surjection $X_U\to U$, effective descent is
\citet[Expose IX, Theorem 4.7]{C_SGA1}. A class in the kernel of pullback
on $H^1$ has a trivialization upstairs; the two pulled-back
trivializations differ by an $M$-translation on
$X_U\times_U X_U$. Its possible descent classes form a subquotient of
$H^0(X_U\times_U X_U,\operatorname{pr}^*M)$. Evaluation at one geometric
point of each connected component bounds this dimension by
$r\#\pi_0(X_U\times_U X_U)$. This proves the required uniform descent
bound, including nonnormal and non-pure-dimensional schemes.

For an open subgroup, Shapiro's lemma replaces $M$ by an induced module
of dimension $[\pi_1(U):H]\dim M$. Finally we use the proved
cohomological crown criterion. A prosolvable group has only abelian
chief factors, so the all-field linear $H^1$ estimate supplies its
finite-generation condition. The argument makes no assertion eliminating
nonabelian crowns of the full group.
\end{proof}
